

\section{Semigroup Proofs of Selected Results}
\label{sec:4.4}

In this final section we present new proofs of several well-known results.
Part one of the next theorem was originally proven by Friedman in \cite{F}. 
To the best of our knowledge, part two of this theorem is a new result.

\begin{theorem}		\label{thm:4.7}
Let $G=(V,\mathcal{E})$ be a strongly connected, $d$-out digraph with adjacency matrix $A$. 
Suppose $w=\left( w_1,w_2,\dots,w_n\right) $ is such that $wA=dw$ and let $W=\sum\limits_{i=1}^nw_i$ be the Friedman weight of $G$. 
We can and do assume that each $w_i$ is a positive integer.

Let $S=\left\langle R_1,R_2,\dots,R_d\right\rangle$ be a coloring semigroup
of $G$ with kernel $\mathcal{K}$ having rank $k$. 
Then 
%
\begin{enumerate}
\item[{\normalfont(1)}] $k\mid W$

\item[{\normalfont(2)}] $k\mid |\mathcal{H}|$ where $\mathcal{H}$ is any maximal group in $\mathcal{K}$.
\end{enumerate}
%
\end{theorem}

\begin{proof}
If $S=\left\langle R_1,R_2,\dots,R_d\right\rangle$ then
\[
  \frac1dA=\frac1d\sum_{j=1}^dR_j.
\]
Notice $\dfrac1dA$ is stochastic, and since strongly connected 
\[
  \lim_{m\to\infty}\frac1m\sum_{k=1}^m\left( \frac1dA\right) ^k=A_0,
\]
where $A_0$ is a rank one strictly positive stochastic matrix with each row
equal to $w_0=\left( \dfrac1W\right) w$.

Now consider the element in the semigroup algebra of binary stochastic
matrices
\[
  \frac1d R_1+\frac1d R_2+\cdots+\frac1d R_d.
\]
Since the coloring semigroup is finite, there is an equivalence between
powers of this element and convolution powers of the probability measure $\mu$ where
\[
  \mu\left( R_i\right) =\frac1d.
\]

It is known that
\[
  \lim_{m\to\infty}\left( \frac1m\sum_{k=1}^m\mu^k\right) =\lambda
\]
where the support of the measure $\lambda$ is precisely the kernel $\mathcal{K}$ of the coloring semigroup $S$. 
(See H\"ogn\"as and Mukherjea \cite{HM}, pages 79 and 92, for more details.) 
Thus
\[
  A_0=\sum_{K\in \mathcal{K}}\lambda(K) K.
\]

Now $\lambda$ is an idempotent probability measure which is well-known to be
locally uniform on each maximal group $\mathcal{H}$.

Precisely, let $E\in E(\mathcal{K})$, then since $A_0$ is a rank one stochastic matrix,
\[
  EA_0E=A_0E.
\]
Also, $\delta_E\ast\lambda\ast\delta_E$, which is an idempotent probability
measure on its support $E\mathcal{K}E$ (the reason being $\lambda\ast\delta_E\ast\lambda=\lambda$, see \cite{HM}), is uniform on $E\mathcal{K}E$; 
thus
\[
  A_0E=\sum_{H\in E\mathcal{K}E}\frac 1{|\mathcal{H}|}\,H. \tag{$\ast$}
\]
Since $\rank(E)=k$, renumber the range elements of $E$ to be $\{1,\dots,k\}$. 
Then
\[
  E=\left[ \left.\dfrac{I_{k\times k}}T\ \right|\ 0\ \right] 
\]
where $T$ is $(b-k)\times k$. 
Also, each $H\in \mathcal{H} = E\mathcal{K}E$ has the same zero columns as $E$, being elements in the same group.

Write $\mathcal{H} = \left\{H_1,H_2,\dots,H_m\right\}$ and let $\mathcal{H}_1=\left\{H_1,H_2,\dots,H_t\right\}$ be all the elements such that
\[
  \left( H_j\right) _{1,1}=1\qquad 1\le j\le t;
\]
then by ($\ast$) above
\[
  \left( A_0E\right) _{1,1}=\frac t{|\mathcal{H}|}.
\]
But $A_0E$ has identical rows, thus
\[
  \left( A_0E\right) _{2,1}=\frac t{|\mathcal{H}|}.
\]

Now $\mathcal{H}_1$ is a subgroup of $\mathcal{H}$ that fixes the element $1$. 
Let $H_{t+1}$ be such that
\[
  \left( H_{t+1}\right) _{1,2}=1
\]
then $\mathcal{H}_1 H_{t+1}=\left\{H_{t+1},\dots,H_{2t}\right\}$ is a coset of $\mathcal{H}_1$.

Notice for $j$ such that $t+1\le j\le 2t$
\[
  \left( H_j^{-1}\right) _{2,1}=1;
\]
thus
\[
  \left( A_0E\right) _{1,2}=\frac t{|\mathcal{H}|}
\]
thus,
\[
  A_0E=\left[ \left.\begin{matrix} 
	\frac t{|\mathcal{H}|} & \frac t{|\mathcal{H}|} 
		&\cdots & \frac t{|\mathcal{H}|} \\
	\vdots & {} & {} & \vdots \\
	\vdots & {} & {} & \vdots \\
	\frac t{|\mathcal{H}|} & \cdots & \cdots & \frac t{|\mathcal{H}|}  
	\end{matrix}\,\right|\ 0\ \right] 
\]
Since $A_0E$ is a stochastic matrix, then
\[
  \frac{tk}{|\mathcal{H}|}=1
\]
which implies
\[
  k\,\bigg|\left|\mathcal{H}\right|.
\]
Also, since
\[
  A_0=\begin{pmatrix} w_0 \\ w_0 \\ \vdots \\ w_0\end{pmatrix}
\]
where
\[
  w_0=\left( \frac{w_1}W,\frac{w_2}W,\cdots,\frac{w_n}W\right) \qquad w_i\in Z^+
\]
then
\[
  A_0E=\left[ \left.\begin{matrix}
	\sum\limits_{E_{i,1}=1}\frac{w_i}W, 
		& \sum\limits_{E_{i,2}=1}\frac{w_i}W, 
		&\cdots, 
		& \sum\limits_{E_{i,k}=1}\frac{w_i}W \\
	\vdots & \vdots & \vdots & \vdots \\
	\vdots & \vdots & \vdots & \vdots \\
	\vdots & \vdots & \vdots & \vdots
	\end{matrix} \right|\,\, 0\right]  
\]

Let
\[
  \sum_{E_{i,1}=1}w_i=W_0.
\]
Then
\[
  \frac{W_0}W=\frac t{|\mathcal{H}|};
\]
but since $|\mathcal{H}|=tk$, then
\[
  kW_0=W
\]
and thus,
\[
  k\mid W. \quad \qed
\]
\end{proof}

In the remainder of the paper, we will make the simplifying assumption that
the out degree of each vertex is two. 
We are leading up to O'Brien's result, in which he assumes that a cycle of prime length exists within the graph. 
Under this assumption, O'Brien also showed that the general problem can be reduced to one of out degree two (see \cite{O} for details).
The cycle $C$ will, for the purposes of notation, always in the rest of this
paper consist of vertices $1\to 2\to 3\cdots\to p\to 1$. 
If $p$ is assumed to be prime, it will be explicitly stated.

\begin{definition}          \label{defn:4.4}
A binary stochastic matrix $R$ will be called a $C$-matrix for $G$ iff
%
\begin{enumerate}
\item[(1)]  $A=R+B$

\item[(2)]  $\left( R\right) _{i,i+1}=\left( R\right) _{p,1}=1$ for $i=1,\dots,p-1$.

\item[(3)]   For each $w\in V$, there is an $i\in C$ such that
$\left( R^t\right) _{w,i}=1$, for some $t$ sufficiently large.
\end{enumerate}
%

It is clear from this definition that for $t_0$ sufficiently large
\[
  R^{t_0}=\left[ \left.\dfrac{I_{p\times p}}M\ \right|\  0\ \right] 
\]
where $I_{p\times p}$ is the $p\times p$ identity matrix and $M$ is a binary
stochastic $(b-p)\times p$ matrix. 
In fact, the powers of $R$ generate an abelian semigroup whose kernel is a cyclic group of order $p$. 
$R^{t_0}$ is simply the identity element for this cyclic group.
\end{definition}

\begin{definition}          \label{defn:4.5}
A vertex $w\in V$ is called periodic with respect to a $C$-matrix $R$ iff
whenever $A_{w,u}=A_{w,v}=1$, then there exists $k\ge 1$ such that for some $i\in C$,
\[
  \left( R^k\right) _{u,i}=\left( R^k\right) _{v,i}=1.
\]
If $w$ is not periodic with respect to $R$, then it is called aperiodic.
\end{definition}

\begin{theorem}		\label{thm:4.8}
If $G$ is aperiodic and $R$ is any $C$-matrix of $G$, then there exists a
vertex $w\in V$ which is aperiodic with respect to $R$.
\end{theorem}

\begin{proof}
Suppose each $w\in V$ is periodic with respect to a given $C$-matrix $R$.
Then if $u$, $v$ are such that
\[
  A_{w,u}=A_{w,v}=1.
\]
Then for some $k$ ($=k(w)$),
\[
  \left( R^k\right) _{u,i}=\left( R^k\right) _{v,i}=1.
\]
Thus
\[
  \left( BR^k\right) _{w,i}=\left( R^{k+1}\right) _{w,i}=1.
\]
But once the $w$-th row of $BR^k$ and $R^{k+1}$ are equal, they remain equal
under additional powers of $R$. 
Since this is true for each vertex, there exists a $\bar{k}=\max\limits_w(k(w))$, such that for all $k\ge\bar{k}$
\[
  BR^k=R^{k+1}.
\]
This implies,
\[
  \begin{gathered}
	\left( \frac12 A\right) R^k =\left( \frac12 R+\frac12 B\right) R^k \\
	=R^{k+1}.
  \end{gathered}
\]

Recall that $A$ is irreducible and aperiodic or, equivalently,
\[
  \left( \frac12 A\right) ^t>0\quad\text{for some }t.
\]
However, since
\[
  \left( \frac12 A\right) ^t R^k=R^{k+t},\quad t\ge 1,
\]
the $p\times p$ submatrix of $R^{k+t}$ (corresponding to the cycle) is
strictly positive. 
Since $R^{k+t}$ is a binary stochastic matrix, we have our contradiction.
\qed
\end{proof}

\begin{theorem}		\label{thm:4.9}
Let $R$ be a $C$-matrix for $G$ and let $\mathcal{K}$ be the kernel of
$S=\left\langle R,B\right\rangle$. Then the rank of 
$\mathcal{K}$ divides $p$. Here $p$ need not be a prime.
\end{theorem}

\begin{proof}
Let 
\[
  M=R^t=\left[ \left.\begin{matrix} 
	0 & 1 & 0 & \cdots & 0 \\
	0 & 0 & 1 & \cdots & \vdots \\
	\vdots & {} & {} & {} & \vdots \\
	0 & \cdots & \cdots & \cdots & 1 \\
	1 & 0 & \cdots & \cdots & 0 \\
	{--} & {--} & {--} & {--} & {--}  \\
	{} & {} & D & {} & {}
	\end{matrix}\,\right|\,0\right] .
\]

Let $z\in \mathcal{K}$ and let $\mathcal{G} = zM\mathcal{K}zM$. 
Suppose $E=E^2$ is the local identity of $\mathcal{G}$. 
Notice that each zero column of $M$ is a zero column of $E$.

Suppose that $\rank(E)=k$ and that $k$ does not divide $p$. 
Then
\[
  p=kr+s,\quad 1\le s<k.
\]

If we can show that, as a result of this assumption, there exists a positive
integer $m$ such that
\[
  \rank\left( EM^mE\right) <k,
\]
then we are done.
\qed
\end{proof}

Let the non-zero columns of $E$ be the columns $i_1<i_2\cdots<i_k\le p$ and
define
\[
  P_j:=\left\{\ell:E_{\ell,i_j}=1,1\le\ell\le p\right\}.
\]
Since the $P_j$'s are disjoint and
\[
  \card\left( \bigcup\limits_{j=1}^k P_j\right) =p<(r+1)k
\]
there must exist at least one $j_0$ such that $\card\left( P_{j_0}\right) \le r$. 
With no loss of generality, assume that $j_0=1$.

\runinhead{Claim.}
For some $m$, $1\le m\le p$,
\[
  \left( M^mE\right) _{i_j,i_1}=0\quad 1\le j\le k.
\]
Notice that, because of the action of the cycle, if $\left( M^mE\right) _{i_j,i_1}=1$,
then $i_j+m$ must be in $P_1$. 
Therefore, $m$ has at most $r$ choices for each $j$. 
Thus, there are at most $rk$ choices and if $m$ is different from each of these choices so that for this $m$, $i_j+m$ is not in $P_1$ for any $j$, $1\le j\le k$
\[
  \left( M^mE\right) _{i_j,i_1}=0.
\]
It follows that the $i_1$th column of $EM^mE$ is zero and this implies that
\[
  \rank\left( EM^mE\right) <k.
\]

\begin{theorem}		\label{thm:4.10}
Suppose $G=(V,\mathcal{E})$ is aperiodic, has no parallel edges and a simple
cycle, $C$, of prime length $p$. 
Then there exists a $C$-matrix $R$ whose coloring semigroup $S=\langle R,B\rangle$ contains a synchronizing instruction.
\end{theorem}

\begin{proof}
Let $A=R+B$. 
Construct the $C$-matrix $R$ in the following way. 
If for $i$ not in $C$, there are $j$ in $C$ and $k$ not in $C$ such that
\[
  A_{i,j}=A_{i,k}=1,
\]
then we choose so that $R_{i,j}=1=B_{i,k}$. 
Now if there are any vertices outside of $C$ not connected to $C$ by the 
previous rule, simply connect them to $C$ by the shortest path to define
$R$.

\runinhead{Case 1.} Every point in $V-C$ is periodic with respect to $R$.

\vphantom{...}

Notice that if every point in $V-C$ is periodic with respect to $R$, then by
the assumption of no parallel edges, and by construction the $(V-C)\times C$ block of the matrix $B$ is all zero. 
(If a vertex in $V-C$ is connected to two distinct vertices in $C$, it must be aperiodic.) 
We know that since every point in $V-C$ is  periodic, $BR^n=R^{n+1}$, for 
sufficiently large $n$, when we restrict the matrices to $V-C$. 
Notice that because of our convention mentioned earlier, for every positive integer $p$ the $(V-C)\times V$ block of $B^{p+1}R^n$ is equal to the product of the $(V-C)\times (V-C)$ block of $B$ and the $(V-C)\times V$ block of $B^pR^n$. 
Thus, by induction, it is easy to see that for any word $H$ in $B$
and $R$ of length $r$,
\[
  \begin{gathered}
	HR^n\text{ restricted to the }(V-C)\times V\text{ block} \\
	=R^{n+r}\text{ restricted to the }(V-C)\times V\text{ block}. 
  \end{gathered}
\]

This means that $[(1/2)A]^p\left( R^n\right) = R^{n+p}$, when we restrict both matrices
to the $(V-C)\times V$ block, for some sufficiently large $n$, but for all positive integers $p$. 
But this is a contradiction as the left side converges as $p$ goes to 
infinity, but the right side does not unless $p=1$. 
The case $p=1$ is trivial to synchronize.

\runinhead{Case 2.} There is at least one aperiodic point outside $C$ with
respect to $R$.

\vphantom{...}

\subruninhead{First observation:} 
Let $w_1$ be aperiodic, $w_1$ not in $C$.
Consider all those vertices $v_k$ outside $C$ such that 
\[
  v_kR^k =w_1.
\]
If there is no such positive $k$, then we work with $w_1$; otherwise, let
$k_0$ be the largest such $k$. 
If $v_{k_0}$ is aperiodic, then we will replace $w_1$ by $v_{k_0}$ and we'll still call it $w_1$. 
If $v_{k_0}$ is periodic, we will not change anything and still work with our old $w_1$. 
Thus, there is no loss of generality in assuming that $w_1$ is aperiodic, $w_1$ not in $C$, and if there is a $v_k$ such that $v_kR^k=w_1$ for positive integer $k$, then such a $v_k$ must be periodic.

\vphantom{...}

\subruninhead{Second observation:} 
Let us consider the shortest path from $C$ to $w_1$. 
We consider a word $H$ (in $R$ and $B$) such that $(C)H=\{w_1,w_2,\dots,w_p\}$, and for $i>1$, the $H$-path from $C$ to $w_i$ does not contain $w_1$. In other words, if we interchange $R$ and $B$ only at the vertex $w_1$, then $(C)H$ remains unaffected.

\vphantom{...}

\subruninhead{Final observation:} 
Decompose $(C)H=\{w_1\}\cup Y_k\cup Z_k$, where $Y_k=\{w_2,\dots,w_k\}$ and
$Z_k=\{w_{k+1},\dots,w_p\}$, such that for $1<j<k+1$, $w_jR^{s_j-1}=w_1$ and
$w_j$ is periodic, where $s_2,s_3,\dots,s_k$ are distinct positive integers, and for each $j>k$, there is no positive integer $p$ such that
$w_jR^p=w_1$. 
[Note: There cannot be two different $w_i$ in $Y_k$ leading to $w_1$ in the same number of $R$-steps, as otherwise we can easily reduce the rank of the kernel.] 
Thus, every vertex in $Y_k$ is periodic; and when we apply the map $R^{n+1}$ on $(C)H$, then the vertices in $Z_k$ remain unaffected if we interchange $R$ and $B$ only at $w_1$.


Suppose $Z_k$ is empty (in this case, recoloring is not needed). 
Choose $n$ large so that $(V)R^n$ is inside $C$ and $w_1R^{n+1}=1$. 
Then, $Y_k$ is mapped onto $\{2,3,\dots,p\}$ by both $R^{n+1}$ and $BR^n$. 
However, $w_1BR^n>1$.

Suppose $Z_k$ is not empty. 
Choose $n$ large as before. 
Let $w_jR^{n+1}=q_j$. Then, $w_1BR^{n+1}=q_{j_0}$, for some $j_0>k$. 
Now we recolor at the vertex $w_1$ and then again consider $(C)HR^{n+1}$, with the new $R$. 
Notice that the image of $Z_k$ remains  the same before and after coloring under the map $R^{n+1}$; however, $q_{j_0}$, the present image of $w_1$ after recoloring, creates a contradiction.
\qed
\end{proof}

Finally, we state a result where the elementary proof is left to the reader.
It is included since it dispenses with a large number of examples quite easily. Notice no assumption concerning parallel edges is required.

\begin{theorem}		\label{thm:4.11}
Let $G=(V,E)$ have a simple cycle $C$ of prime length $p$ and adjacency
matrix $A$. 
Suppose one of the following conditions holds:
%
\begin{enumerate}
\item[{\normalfont(1)}]  there exists $j\notin C$ such that
\[
  \sum_{i=1}^p A_{i,j}\ge 2\text{; or}
\]

\item[{\normalfont(2)}]  there exists $j\in C$ such that
\[
  \sum_{i=1}^p A_{i,j}\ge 3.
\]
\end{enumerate}
%
Then $G$ has a coloring semigroup that contains a synchronizing instruction.
\end{theorem}

\paragraph{Acknowledgements}

The authors are grateful to the members of the Discrete Mathematics Seminar at the University of South Florida (USF) for their kind attention to a series of lectures given 
by the authors.  
In particular, conversations with W.E.\ Clark, N.\ Jonaska, and S.\ Suen were always interesting.  
The authors would  also like to thank G\"oran H\"ogn\"as for helpful comments.

This work was done while the first author was visiting USF while on sabbatical from Southern Illinois University at Carbondale (SIUC).
Both the generosity of SIUC and the hospitality of USF were very much appreciated. 
